% ============================================================================
% Fichier  : partie4/P04_C22_operations_incidents.tex
% Titre    : Opérations de Sécurité et Incidents
% Auteur   : MINKA MI NGUIDJOÏ Thierry Emmanuel
% Version  : 1.0 -- novembre 2026
% ============================================================================
% Base : brouillon « Livre VIII » (opérations de sécurité). Sources : NIST SP
% 800-61 (Rev. 2, 2012 ; Rev. 3, 2025, alignée sur le CSF 2.0) ; ISO/IEC
% 27035-1:2023 ; NIST SP 800-92. Délai de détection de l'Opération GAROUA
% calculé d'après la chronologie du Prologue : 80 h 18 min.
% ============================================================================

\chapter{Opérations de Sécurité et Incidents}
\label{chap:operations}
\chaptermark{Opérations et incidents}

\epigraph{%
    « On ne choisit pas d'être attaqué~;\\
    on choisit d'être prêt. »%
}{Principe du cours}

\niveauChapitre{Intermédiaire}{Chapitres~\ref{chap:menaces},
\ref{chap:reseaux2} et~\ref{chap:systemes}.}

\begin{boxobjectifs}
\begin{itemize}
    \item Décrire l'organisation et les outils d'un centre opérationnel de
          sécurité.
    \item Écrire une règle de détection à partir d'une technique
          d'attaque.
    \item Dérouler le cycle de réponse aux incidents du NIST et de
          l'ISO/IEC~27035.
    \item Classer un incident et décider d'une escalade.
    \item Rédiger un guide de réponse à un rançongiciel.
    \item Mesurer la performance de la détection et de la réponse.
\end{itemize}
\end{boxobjectifs}

\bigskip

% ============================================================================
\section*{Question directrice : pourquoi a-t-il fallu trois jours pour voir l'attaquant~?}
% ============================================================================

L'attaquant s'est connecté le 2~mars à 23~h~12~; LiZbiZ l'a découvert le
6~mars à 7~h~30, soit 80~heures et 18~minutes plus tard, et seulement
parce que tout était chiffré. Pendant ces trois jours, les journaux
existaient sur certaines machines, mais personne ne les lisait. Un
dispositif d'opérations de sécurité transforme des journaux en alertes, et
des alertes en décisions.

\medskip
\textbf{Plan du chapitre.}
\begin{enumerate}
    \item Le centre opérationnel de sécurité (section~\ref{sec:c22-soc}).
    \item Collecter et corréler~: le SIEM (section~\ref{sec:c22-siem}).
    \item Le cycle de réponse aux incidents (section~\ref{sec:c22-cycle}).
    \item Classer et escalader (section~\ref{sec:c22-classer}).
    \item Le guide de réponse au rançongiciel (section~\ref{sec:c22-guide}).
    \item Mesurer (section~\ref{sec:c22-mesurer}).
\end{enumerate}

\separateur

% ============================================================================
\section{Le Centre Opérationnel de Sécurité}
\label{sec:c22-soc}
% ============================================================================

\begin{boxdef}
\termecle{Centre opérationnel de sécurité} (SOC)~: équipe, processus et
outils qui surveillent en continu le système d'information, détectent les
événements de sécurité, les analysent et déclenchent la réponse.
\end{boxdef}

\begin{tableaucours}{L{2.6cm}L{9.6cm}}{%
    \tableauHeader{Niveau} &
    \tableauHeader{Rôle}
}
Niveau~1 & Tri des alertes~: vrai ou faux positif, première
qualification, escalade \\
Niveau~2 & Analyse approfondie, corrélation, confinement initial \\
Niveau~3 & Recherche proactive de menaces, investigation, réglage des
règles de détection \\
Responsable & Organisation, indicateurs, liens avec la direction et le
RSSI \\
\end{tableaucours}
\vspace{-0.8em}
\begin{center}{\small\itshape Tableau~22.1. L'organisation d'un SOC}\end{center}

Une entreprise de la taille de LiZbiZ n'a pas les moyens d'un SOC interne
fonctionnant jour et nuit~; elle peut en confier la surveillance à un
prestataire, à condition de garder la maîtrise des décisions et des
journaux.

% ============================================================================
\section{Collecter et Corréler : le SIEM}
\label{sec:c22-siem}
% ============================================================================

Un \termecle{SIEM} (gestion des informations et des événements de
sécurité) collecte les journaux de toutes les sources, les met dans un
format commun, les conserve, et applique des règles de corrélation qui
produisent des alertes~\cite{nist80092}.

\begin{boxoutils}{Une règle de détection, exprimée en pseudo-code}
\begin{lstlisting}
regle "Connexion distante hors horaires suivie d'un nouvel administrateur"
  si   evenement = connexion_vpn_reussie
   et  heure hors de [07:00 ; 19:00]
   et  dans les 72 h, sur le meme compte ou depuis la meme adresse :
       evenement = ajout_au_groupe "Admins du domaine"
  alors alerte de severite haute, niveau 2
\end{lstlisting}
Cette règle couvre deux tactiques ATT\&CK de l'Opération GAROUA~: l'accès
initial (connexion de 23~h~12) et l'élévation de privilèges (le 4~mars).
Elle n'aurait fonctionné que si la passerelle et les contrôleurs de domaine
envoyaient leurs journaux au SIEM.
\end{boxoutils}

% ============================================================================
\section{Le Cycle de Réponse aux Incidents}
\label{sec:c22-cycle}
% ============================================================================

Le NIST décrit la réponse aux incidents en quatre phases~; sa révision de
2025 l'aligne sur les six fonctions du CSF~2.0~\cite{nist80061r2,nist80061r3}.
L'ISO/IEC~27035 en donne une version normalisée~\cite{iso27035_2023}.

\begin{tableaucours}{L{3.8cm}L{8.4cm}}{%
    \tableauHeader{Phase} &
    \tableauHeader{Contenu}
}
Préparation & Politique, équipe, contacts, outils, guides, exercices \\
Détection et analyse & Repérer, qualifier, établir le périmètre,
documenter \\
Confinement, éradication, rétablissement & Isoler, supprimer la cause,
restaurer à partir de sources saines, surveiller \\
Activité post-incident & Retour d'expérience, amélioration des mesures et
des guides \\
\end{tableaucours}
\vspace{-0.8em}
\begin{center}{\small\itshape Tableau~22.2. Le cycle de réponse aux incidents
(NIST SP~800-61)}\end{center}

\begin{boxalert}
\textbf{Violation typique.} Éradiquer avant d'avoir compris (formater et
réinstaller la première machine infectée).

\smallskip
\textbf{Règle~:} on confine d'abord (isoler du réseau, sans éteindre),
on préserve les preuves (Chapitre~\ref{chap:investigation}), on établit
le périmètre complet, puis on éradique partout à la fois. Sinon,
l'attaquant revient par la machine oubliée.
\end{boxalert}

% ============================================================================
\section{Classer et Escalader}
\label{sec:c22-classer}
% ============================================================================

\begin{tableaucours}{L{2.0cm}L{6.4cm}L{3.8cm}}{%
    \tableauHeader{Gravité} &
    \tableauHeader{Critère (exemple)} &
    \tableauHeader{Escalade}
}
Critique & Activité arrêtée, données sensibles exfiltrées, plusieurs sites
& Direction générale, cellule de crise, sous 1~heure \\
Haute & Compte à privilèges compromis, serveur touché & RSSI et DSI, sous
4~heures \\
Moyenne & Poste infecté isolé, sans propagation & Équipe informatique,
dans la journée \\
Basse & Tentative bloquée, hameçonnage signalé & Traitement courant \\
\end{tableaucours}
\vspace{-0.8em}
\begin{center}{\small\itshape Tableau~22.3. Grille de gravité des incidents
(convention à adapter)}\end{center}

% ============================================================================
\section{Le Guide de Réponse au Rançongiciel}
\label{sec:c22-guide}
% ============================================================================

\begin{boxprocedure}
\textbf{Rançongiciel~: les premières heures}
\begin{enumerate}
    \item \textbf{Isoler} les machines touchées du réseau (débrancher ou
          isoler par l'EDR), \emph{sans les éteindre}.
    \item \textbf{Couper} les accès distants et les liaisons entre sites
          pour stopper la propagation.
    \item \textbf{Protéger les sauvegardes}~: les déconnecter, vérifier
          qu'elles sont intactes.
    \item \textbf{Prévenir} le RSSI et la direction selon la grille~;
          activer la cellule de crise si l'incident est critique.
    \item \textbf{Préserver} les preuves~: mémoire, journaux, note de rançon,
          avec leurs empreintes.
    \item \textbf{Établir le périmètre}~: comptes utilisés, machines
          touchées, données sorties.
    \item \textbf{Réinitialiser} les mots de passe des comptes à privilèges
          depuis un poste sain.
    \item \textbf{Restaurer} à partir de sauvegardes saines, après
          éradication, et surveiller.
\end{enumerate}
\end{boxprocedure}

\begin{boxscenario}{Opération GAROUA. Ce qu'un guide aurait changé}
\textbf{Constatations.} À 22~h~05, le technicien d'astreinte redémarre les
serveurs~; à 7~h~30, les liaisons entre sites sont toujours ouvertes~; les
mots de passe d'administration ne sont changés que le 9~mars.

\medskip
\textbf{Analyse} (confiance moyenne). Avec le guide, l'isolement aurait
remplacé le redémarrage (mémoire préservée), la coupure des liaisons aurait
protégé Yaoundé et Bafoussam, et la réinitialisation des comptes aurait
empêché un retour de l'attaquant.
\end{boxscenario}

% ============================================================================
\section{Mesurer}
\label{sec:c22-mesurer}
% ============================================================================

\begin{tableaucours}{L{4.6cm}L{7.6cm}}{%
    \tableauHeader{Indicateur} &
    \tableauHeader{Définition}
}
Délai de détection (MTTD) & Durée moyenne entre le début d'un incident et
sa détection \\
Délai de réponse (MTTR) & Durée moyenne entre la détection et le
rétablissement \\
Taux de faux positifs & Part des alertes qui ne correspondaient à aucun
incident \\
Couverture des journaux & Part des systèmes critiques qui envoient leurs
journaux au SIEM \\
Couverture des techniques & Part des techniques ATT\&CK prioritaires
couvertes par une règle \\
\end{tableaucours}
\vspace{-0.8em}
\begin{center}{\small\itshape Tableau~22.4. Indicateurs des opérations de sécurité}\end{center}

\begin{boxmetier}{Analyste SOC}
\textbf{Documenter chaque alerte.} L'analyste consigne pour chaque alerte
sa décision (faux positif, incident), ses raisons et ses sources~; ces
traces servent au réglage des règles et à l'enquête.
\end{boxmetier}

\begin{boxmetier}{RSSI}
\textbf{Préparer.} Le RSSI fait écrire les guides, tient à jour la liste
des contacts (direction, juristes, prestataires, autorités), et organise
chaque année un exercice sur table.
\end{boxmetier}

\separateur

\begin{boxresume}
\textbf{Ce qu'il faut retenir de ce chapitre~:}
\begin{itemize}
    \item \textbf{Un SOC} trie, analyse et répond, en niveaux~; il peut être
          externalisé, pas la décision.
    \item \textbf{Le SIEM} centralise et corrèle~; une règle vaut ce que
          valent les journaux qu'il reçoit.
    \item \textbf{La réponse} suit quatre phases~: préparer, détecter et
          analyser, confiner-éradiquer-rétablir, tirer les leçons.
    \item \textbf{On confine avant d'éradiquer}, et l'on préserve les
          preuves.
    \item \textbf{Un guide de réponse} fixe les gestes des premières heures.
    \item \textbf{Délai de détection et délai de réponse} mesurent
          l'efficacité.
\end{itemize}
\end{boxresume}

\begin{boxinfo}
\textbf{Points de contrôle de l'élève.} Je sais~: décrire les niveaux d'un
SOC~; écrire une règle de corrélation~; nommer les phases de la réponse~;
classer un incident~; dérouler le guide rançongiciel~; calculer un délai
de détection.
\end{boxinfo}

\bigskip

\begin{boxexo}[title={\ding{45}\quad Exercice~22.1. Le délai \niveaubase}]
Recalculez le délai de détection de l'Opération GAROUA à partir de la
chronologie du Prologue.
\end{boxexo}

\begin{boxexo}[title={\ding{45}\quad Exercice~22.2. Les phases \niveaubase}]
Rangez dans les phases du NIST~: réinitialiser les mots de passe~; écrire
le guide rançongiciel~; isoler un serveur~; rédiger le retour
d'expérience~; qualifier une alerte.
\end{boxexo}

\begin{boxexo}[title={\ding{45}\quad Exercice~22.3. Classer \niveaubase}]
Classez avec la grille du Tableau~22.3~: un courriel d'hameçonnage
signalé~; un compte d'administrateur utilisé depuis l'étranger~; un poste
de l'agence de Bafoussam infecté par un virus isolé.
\end{boxexo}

\begin{boxexo}[title={\ding{45}\quad Exercice~22.4. Une règle \niveaumoyen}]
Écrivez une règle de détection pour l'exfiltration~: un volume sortant
inhabituel depuis un serveur de fichiers.
\end{boxexo}

\begin{boxexo}[title={\ding{45}\quad Exercice~22.5. Les journaux nécessaires \niveaumoyen}]
Pour chaque règle de ce chapitre, listez les sources de journaux
nécessaires et vérifiez qu'elles existaient chez LiZbiZ en mars.
\end{boxexo}

\begin{boxexo}[title={\ding{45}\quad Exercice~22.6. Confiner ou éradiquer \niveaumoyen}]
Un analyste propose de réinstaller immédiatement le poste infecté qu'il
vient de trouver. Que lui répondez-vous~?
\end{boxexo}

\begin{boxexo}[title={\ding{45}\quad Exercice~22.7. Externaliser \niveaumoyen}]
LiZbiZ veut confier sa surveillance à un SOC externe. Quelles clauses et
quels indicateurs exiger~?
\end{boxexo}

\begin{boxexo}[title={\ding{45}\quad Exercice~22.8. Faux positifs \niveaumoyen}]
Une règle produit 200~alertes par semaine, dont 3~vraies. Proposez trois
façons de la régler sans perdre les vraies alertes.
\end{boxexo}

\begin{boxexo}[title={\ding{45}\quad Exercice~22.9. Un guide \niveauexpert}]
Rédigez le guide de réponse à la compromission d'une boîte de
messagerie (règle de transfert suspecte, Chapitre~\ref{chap:nuage}).
\end{boxexo}

\begin{boxexo}[title={\ding{45}\quad Exercice~22.10. L'exercice sur table \niveauexpert}]
Concevez un exercice sur table de deux heures pour la direction de LiZbiZ,
fondé sur l'Opération GAROUA~: injections, rôles, décisions attendues,
critères d'évaluation.
\end{boxexo}

\bigskip

\begin{boxinfo}
\textbf{Pour aller plus loin.}
\begin{itemize}
    \item NIST SP~800-61 (révisions 2 et~3)~\cite{nist80061r2,nist80061r3}.
    \item ISO/IEC~27035-1~\cite{iso27035_2023}.
    \item NIST SP~800-92, gestion des journaux~\cite{nist80092}.
\end{itemize}
\end{boxinfo}

\begin{boxinfo}
\textbf{Transition.} Après l'incident vient la reprise. La
Partie~\ref{part:avance} commence par la continuité et la reprise
d'activité (Chapitre~\ref{chap:continuite}).
\end{boxinfo}
